\begin{table}[htbp]
\centering
\caption{Lagged Hiring Test: Prior-Year Outcomes by Future Donor Status}
\label{tab:donor_lagged}
\small
\begin{threeparttable}

\textbf{Panel A: Unconditional Comparison}
\medskip

\begin{tabular}{lcccccc}
\toprule
 & \multicolumn{2}{c}{Future Donor} & \multicolumn{2}{c}{Non-Donor} & & \\
 & Mean & SD & Mean & SD & Diff & $p$-value \\
\midrule
Has New Hire (=1) & 0.696 & 0.460 & 0.478 & 0.500 & 0.218 & $<$0.001 \\
No. New Hires & 7.447 & 17.243 & 1.867 & 4.813 & 5.580 & $<$0.001 \\
Avg Tenure (years) & 2.724 & 1.294 & 3.237 & 1.908 & -0.513 & $<$0.001 \\
Female Share & 0.034 & 0.098 & 0.046 & 0.158 & -0.012 & $<$0.001 \\
\bottomrule
\end{tabular}

\bigskip

\textbf{Panel B: Conditional Comparison (Kyoku + Position + Year FE + Log Headcount at $t$)}
\medskip

\begin{tabular}{lccccc}
\toprule
 & Coef & SE & $t$-stat & $p$-value & $N$ \\
\midrule
Has New Hire (=1) & 0.1603$^{***}$ & 0.0156 & 10.29 & $<$0.001 & 7,235 \\
No. New Hires & 5.1419$^{***}$ & 0.6665 & 7.71 & $<$0.001 & 7,235 \\
Avg Tenure (years) & 0.1133$^{***}$ & 0.0289 & 3.92 & $<$0.001 & 7,235 \\
Female Share & -0.0081$^{**}$ & 0.0032 & -2.53 & 0.012 & 7,235 \\
\bottomrule
\end{tabular}

\begin{tablenotes}[flushleft]
\footnotesize
\item \textit{Notes:} This table tests whether the new-hiring gap between donor and non-donor offices reflects vacancy chains or selection. Outcomes (new hires, tenure, female share) are measured at year $t{-}1$; donor status is defined at year $t$ (whether the office had any transfers out). If offices that will become donors already exhibit higher hiring rates in the prior year, this suggests selection rather than vacancy chains. Sample: ka $\times$ position, outcome years 1937--1944, donor status years 1938--1945. Panel~A reports unconditional means. Panel~B reports coefficients from OLS regressions with kyoku, position, and year fixed effects, controlling for log office headcount at year $t$. Standard errors clustered at the office level. $^{***}p<0.01$, $^{**}p<0.05$, $^{*}p<0.1$.
\end{tablenotes}
\end{threeparttable}
\end{table}
